\documentclass[11pt,a4paper]{article}
\usepackage[utf8]{inputenc}
\usepackage[vietnamese]{babel}
\usepackage{amsmath,amssymb,amsfonts,amsthm}
\usepackage{geometry}
\geometry{top=2cm, bottom=2cm, left=2cm, right=2cm}
\usepackage{hyperref}
\usepackage{xcolor}
\usepackage{booktabs}
\usepackage{tcolorbox}
\usepackage{fancyhdr}
\usepackage{array}

% ── ĐỊNH NGHĨA MÀU SẮC CHUYÊN NGHIỆP ──
\definecolor{primary}{HTML}{1A365D}   % Navy đậm
\definecolor{secondary}{HTML}{2B6CB0} % Xanh lam
\definecolor{accent}{HTML}{C53030}    % Đỏ đô điểm nhấn
\definecolor{bglight}{HTML}{F7FAFC}   % Nền xám nhạt
\definecolor{boxborder}{HTML}{CBD5E0} % Viền hộp

\hypersetup{
    colorlinks=true,
    linkcolor=secondary,
    citecolor=secondary,
    urlcolor=secondary
}

% ── ĐỊNH DẠNG TRANG & HEADER/FOOTER ──
\pagestyle{fancy}
\fancyhf{}
\fancyhead[L]{\small \textbf{CS115: Toán cho Khoa học Máy tính}}
\fancyhead[R]{\small Tài liệu Tổng hợp Toàn diện 4 Trụ cột}
\fancyfoot[C]{\thepage}
\fancyfoot[R]{\small \textcolor{gray}{UIT - VNUHCM}}

% ── HỘP THÔNG TIN (TCOLORBOX) ──
\tcbuselibrary{skins,breakable}

\newtcolorbox{mindsetbox}[1][]{
    colback=bglight,
    colframe=secondary,
    fonttitle=\bfseries\color{white},
    title={🧠 AI Mindset \& Tư duy Cốt lõi: #1},
    enhanced,
    attach boxed title to top left={yshift=-2mm, xshift=3mm},
    boxed title style={colback=secondary,colframe=secondary},
    top=4mm, bottom=3mm, left=3mm, right=3mm,
    breakable
}

\newtcolorbox{bottleneckbox}[1][]{
    colback=white,
    colframe=accent,
    fonttitle=\bfseries\color{white},
    title={⚠️ Nút thắt Cốt lõi \& Lỗi thường gặp: #1},
    enhanced,
    attach boxed title to top left={yshift=-2mm, xshift=3mm},
    boxed title style={colback=accent,colframe=accent},
    top=4mm, bottom=3mm, left=3mm, right=3mm,
    breakable
}

\newtcolorbox{symbolbox}[1][]{
    colback=bglight,
    colframe=primary,
    fonttitle=\bfseries\color{white},
    title={📖 Giải mã Ký hiệu Toán học: #1},
    enhanced,
    attach boxed title to top left={yshift=-2mm, xshift=3mm},
    boxed title style={colback=primary,colframe=primary},
    top=4mm, bottom=3mm, left=3mm, right=3mm,
    breakable
}

\newtcolorbox{examplebox}[1][]{
    colback=bglight,
    colframe=primary!80!black,
    fonttitle=\bfseries\color{white},
    title={🔢 Ví dụ Minh họa Thực tế bằng Số: #1},
    enhanced,
    attach boxed title to top left={yshift=-2mm, xshift=3mm},
    boxed title style={colback=primary!80!black,colframe=primary!80!black},
    top=4mm, bottom=3mm, left=3mm, right=3mm,
    breakable
}

% ── TIÊU ĐỀ TÀI LIỆU ──
\title{\Huge \textbf{CS115: TOÁN CHO KHOA HỌC MÁY TÍNH}\\[0.3em]
\Large Cẩm nang Tổng hợp Chuyên sâu 4 Trụ cột Kiến thức, Từ điển Ký hiệu, Chuỗi Logic \& Phương pháp Giảng dạy}
\author{\textbf{Chương trình Cử nhân Khoa học Máy tính / Trí tuệ Nhân tạo}\\Đại học Công nghệ Thông tin (UIT - VNUHCM)}
\date{\today}

\begin{document}

\maketitle

\begin{abstract}
Tài liệu này là công trình tổng hợp hệ thống và toàn diện dựa trên 18 nguồn tài liệu bài giảng gốc của môn học \textbf{CS115 -- Toán cho Khoa học Máy tính}. Được thiết kế đặc biệt nhằm tháo gỡ các nút thắt hiểu biết (contract bottlenecks) và các lỗi tư duy thường gặp cho sinh viên mới, tài liệu phân tích chi tiết 4 trụ cột toán học cốt lõi: (1) Đại số Tuyến tính \& Giảm chiều Dữ liệu, (2) Giải tích Vector, Ma trận Hessian/Jacobian \& Tối ưu hóa, (3) Xác suất Thống kê, Phân phối \& Ước lượng Tham số (MLE/MAP/Bayes), và (4) Mô hình Học máy, Hồi quy, Phân loại \& Mạng Neural (Backpropagation). Mỗi khái niệm đều đi kèm bộ từ điển ký hiệu, chuỗi logic liên hoàn (Logic Chain), tư duy kỹ sư AI (AI Mindset) và các ví dụ tính toán số từng bước trực quan.
\end{abstract}

\tableofcontents
\newpage

% ════════════════════════════════════════════════════════════════
\section*{Lời nói đầu: Khai phóng "Nút thắt Cốt lõi" của Môn học}
\addcontentsline{toc}{section}{Lời nói đầu: Khai phóng "Nút thắt Cốt lõi" của Môn học}

Trong kỷ nguyên Trí tuệ Nhân tạo (AI), môn học \textbf{CS115 -- Toán cho Khoa học Máy tính} đóng vai trò là "bộ não vận hành" đằng sau mọi thuật toán Học máy (Machine Learning - ML) và Học sâu (Deep Learning - DL). Rất nhiều sinh viên mới tiếp xúc thường gặp e ngại trước vô số ký hiệu toán học rời rạc. Bản chất sai lầm không nằm ở khả năng tính toán, mà ở việc thiếu một \textbf{Bức tranh Tổng thể (Big Picture)} kết nối các mảng kiến thức.

\begin{mindsetbox}[Sơ đồ Tư duy 4 Bước của Trí tuệ Nhân tạo]
Một mô hình AI tự học luôn vận hành theo chuỗi logic 4 mắt xích liên hoàn:
$$\text{Dữ liệu thô } \mathbf{x} \xrightarrow{\text{(1) Đại số}} \text{Vector/Ma trận} \xrightarrow{\text{(2) Thống kê}} \text{Dự đoán xác suất } \hat{y} \xrightarrow{\text{(3) Hàm Loss}} \text{Mức độ sai hỏng } L \xrightarrow{\text{(4) Giải tích/Tối ưu}} \text{Sửa trọng số } \mathbf{w}$$
\begin{enumerate}
    \item \textbf{Biểu diễn (Representation - Đại số tuyến tính):} Chuyển văn bản, ảnh, âm thanh thành các Vector/Ma trận số.
    \item \textbf{Dự đoán (Inference - Xác suất & Hồi quy):} Mô hình tính điểm số và chuyển đổi thành xác suất tin cậy.
    \item \textbf{Đo lường (Evaluation - Loss Function):} So sánh dự đoán $\hat{y}$ với thực tế $y$ để tính sai số $L$.
    \item \textbf{Tự học (Optimization - Vi tích phân & Gradient Descent):} Tìm hướng dốc nhất để điều chỉnh các trọng số $\mathbf{w}$ làm sai số giảm dần về 0.
\end{enumerate}
\end{mindsetbox}

% ════════════════════════════════════════════════════════════════
\section{Bảng Rà soát Nút thắt hiểu biết \& Các Lỗi thường gặp (Pitfalls)}

Dưới đây là 6 nút thắt tư duy và lỗi toán học kinh điển nhất được rà soát đối soát từ dữ liệu bài thi và thực hành CS115:

\begin{bottleneckbox}[6 Lỗi Toán học Kỹ thuật Sinh viên Rất Dễ Mắc Phải]
\begin{enumerate}
    \item \textbf{Quên chuẩn hóa Vector trong Tiến trình Gram-Schmidt:}
    \textit{Sai lầm:} Sinh viên chỉ tính triệt tiêu bóng chiếu để ra hệ trực giao $\{\mathbf{v}_1, \mathbf{v}_2, \dots, \mathbf{v}_p\}$ rồi dừng lại.\\
    \textit{Xử lý chuẩn xác:} Phải thực hiện bước chia cho độ dài $\mathbf{u}_i = \frac{\mathbf{v}_i}{\|\mathbf{v}_i\|_2}$ để đưa về cơ sở \textbf{Trực chuẩn (Orthonormal)}.
    
    \item \textbf{Nhầm lẫn giữa Ma trận Jacobian và Ma trận Hessian:}
    \textit{Sai lầm:} Xem Jacobian là đạo hàm cấp 2 hoặc Hessian là ma trận đạo hàm của hàm vector.\\
    \textit{Xử lý chuẩn xác:} Ma trận \textbf{Jacobian ($\mathbf{J}$)} là bảng đạo hàm \textbf{cấp 1} của hàm \textbf{vector} ($\mathbb{R}^D \rightarrow \mathbb{R}^E$). Ma trận \textbf{Hessian ($\mathbf{H}_f$)} là bảng đạo hàm \textbf{cấp 2} của hàm \textbf{vô hướng} ($\mathbb{R}^D \rightarrow \mathbb{R}$).
    
    \item \textbf{Dùng nhầm Hàm Mất Mát (MSE) cho Bài toán Phân loại:}
    \textit{Sai lầm:} Áp dụng MSE Loss $(\hat{y} - y)^2$ cho Logistic Regression hoặc Softmax Classifier.\\
    \textit{Xử lý chuẩn xác:} Hàm Sigmoid/Softmax phi tuyến kết hợp MSE tạo ra địa hình sai số không lồi (non-convex) có nhiều hố sâu địa phương. Phải sử dụng \textbf{Cross-Entropy Loss (Negative Log-Likelihood)} để đảm bảo địa hình lồi, giúp Gradient Descent hội tụ toàn cục.
    
    \item \textbf{Bỏ qua Dấu của Định thức con khi Xét Dấu Ma trận (Sylvester's Criterion):}
    \textit{Sai lầm:} Chỉ tính định thức toàn phần $\det(\mathbf{H}_f)$ để kết luận cực trị.\\
    \textit{Xử lý chuẩn xác:} Để ma trận Hessian \textbf{Xác định dương} (Cực tiểu), \textbf{TẤT CẢ} định thức con góc trên bên trái phải dương ($\Delta_1 > 0, \Delta_2 > 0, \Delta_3 > 0$). Để \textbf{Xác định âm} (Cực đại), chuỗi định thức con phải \textbf{đan dấu} ($\Delta_1 < 0, \Delta_2 > 0, \Delta_3 < 0$).
    
    \item \textbf{Nhầm lẫn giữa PCA và SVD:}
    \textit{Sai lầm:} Coi PCA và SVD là hai thuật toán độc lập hoàn toàn.\\
    \textit{Xử lý chuẩn xác:} \textbf{SVD} là phép phân rã ma trận tổng quát $\mathbf{A} = \mathbf{U} \mathbf{\Sigma} \mathbf{V}^T$. \textbf{PCA} thực chất chính là phép SVD áp dụng trên ma trận dữ liệu đã được chuẩn hóa trừ trung bình (Zero-mean data).
    
    \item \textbf{Quên nhân Đạo hàm Địa phương với Upstream Gradient trong Backpropagation:}
    \textit{Sai lầm:} Khi truyền lùi qua các cổng, sinh viên chỉ tính đạo hàm riêng của cổng mà quên nhân với tín hiệu lỗi từ phía sau tới.\\
    \textit{Xử lý chuẩn xác:} Theo Chain Rule: $\text{Downstream Grad} = \text{Upstream Grad} \times \text{Local Grad}$.
\end{enumerate}
\end{bottleneckbox}

\newpage
% ════════════════════════════════════════════════════════════════
\section{Trụ cột 1: Đại số Tuyến tính, Phép chiếu \& Giảm chiều Dữ liệu}

\subsection{Từ điển Ký hiệu Đại số Tuyến tính}
\begin{symbolbox}[Các Ký hiệu Biểu diễn Dữ liệu]
\begin{itemize}
    \item $\mathbb{R}$: Tập số thực. $\mathbb{R}^D$: Không gian vector $D$ chiều (danh sách $D$ số thực).
    \item $\mathbf{x} \in \mathbb{R}^D$: Vector cột (viết chữ in đậm thường). $\mathbf{x}^T = [x_1, x_2, \dots, x_D]$ là vector hàng chuyển vị.
    \item $\mathbf{A} \in \mathbb{R}^{M \times N}$: Ma trận kích thước $M$ hàng và $N$ cột (viết chữ in đậm hoa).
    \item $\mathbf{u}^T \mathbf{v} = \sum_{i=1}^D u_i v_i$: Tích vô hướng (Dot product). Đo độ tương đồng về hướng giữa 2 vector.
    \item $\|\mathbf{x}\|_2 = \sqrt{\sum_{i=1}^D x_i^2}$: Chuẩn $L_2$ (Khoảng cách Euclidean đường chim bay).
    \item $\|\mathbf{x}\|_1 = \sum_{i=1}^D |x_i|$: Chuẩn $L_1$ (Khoảng cách Manhattan ô bàn cờ).
    \item $\mathbf{I}_N$: Ma trận đơn vị kích thước $N \times N$ (đường chéo chính bằng 1, các vị trí khác bằng 0).
\end{itemize}
\end{symbolbox}

\subsection{Phép chiếu Trực giao (Orthogonal Projection)}
\textbf{Khái niệm & Mindset:} Khi một vector $\mathbf{y}$ không nằm trong không gian con $W$, vị trí đại diện tốt nhất cho $\mathbf{y}$ trong $W$ chính là bóng chiếu vuông góc $\text{proj}_W \mathbf{y}$.

\textbf{Công thức tính toán:}
Hình chiếu của vector $\mathbf{y}$ lên đường thẳng chứa vector $\mathbf{u} \neq \mathbf{0}$:
$$\text{proj}_{\mathbf{u}} \mathbf{y} = \frac{\langle \mathbf{y}, \mathbf{u} \rangle}{\|\mathbf{u}\|_2^2} \mathbf{u} = \frac{\mathbf{y}^T \mathbf{u}}{\mathbf{u}^T \mathbf{u}} \mathbf{u}$$
Vector độ lệch $\mathbf{z} = \mathbf{y} - \text{proj}_{\mathbf{u}} \mathbf{y}$ hoàn toàn vuông góc với $\mathbf{u}$ (tức $\mathbf{z}^T \mathbf{u} = 0$).

\begin{examplebox}[Tính Phép chiếu Trực giao]
Cho $\mathbf{y} = (6, 2, 4)^T$ và $\mathbf{u} = (1, 2, 0)^T$.
\begin{enumerate}
    \item Tích vô hướng: $\mathbf{y}^T \mathbf{u} = (6 \times 1) + (2 \times 2) + (4 \times 0) = 10$.
    \item Bình phương độ dài: $\mathbf{u}^T \mathbf{u} = 1^2 + 2^2 + 0^2 = 5$.
    \item Hình chiếu: $\text{proj}_{\mathbf{u}} \mathbf{y} = \frac{10}{5} (1, 2, 0)^T = 2(1, 2, 0)^T = (2, 4, 0)^T$.
    \item Kiểm tra trực giao: Vector độ lệch $\mathbf{z} = (6-2, 2-4, 4-0)^T = (4, -2, 4)^T$. Tích $\mathbf{z}^T \mathbf{u} = (4 \times 1) + (-2 \times 2) + (4 \times 0) = 0$ (Chính xác vuông góc!).
\end{enumerate}
\end{examplebox}

\subsection{Tiến trình Gram-Schmidt (Gram-Schmidt Process)}
\textbf{Chuỗi Logic:} Chuyển một cơ sở xiên xẹo $\{\mathbf{x}_1, \mathbf{x}_2, \dots, \mathbf{x}_p\}$ thành cơ sở trực giao $\{\mathbf{v}_1, \mathbf{v}_2, \dots, \mathbf{v}_p\}$, sau đó chia độ dài để ra cơ sở trực chuẩn $\{\mathbf{u}_1, \mathbf{u}_2, \dots, \mathbf{u}_p\}$.

\textbf{Thuật toán từng bước:}
\begin{align*}
\mathbf{v}_1 &= \mathbf{x}_1 \\
\mathbf{v}_2 &= \mathbf{x}_2 - \frac{\langle \mathbf{x}_2, \mathbf{v}_1 \rangle}{\|\mathbf{v}_1\|^2} \mathbf{v}_1 \\
\mathbf{v}_k &= \mathbf{x}_k - \sum_{j=1}^{k-1} \frac{\langle \mathbf{x}_k, \mathbf{v}_j \rangle}{\|\mathbf{v}_j\|^2} \mathbf{v}_j \quad (k = 2, \dots, p) \\
\mathbf{u}_i &= \frac{\mathbf{v}_i}{\|\mathbf{v}_i\|_2} \quad (\text{Bước Trực chuẩn hóa})
\end{align*}

\subsection{Phân rã SVD \& Phân tích Thành phần Chính PCA}
\textbf{SVD (Singular Value Decomposition):} Phân rã ma trận bất kỳ $\mathbf{A} \in \mathbb{R}^{M \times N}$:
$$\mathbf{A} = \mathbf{U} \mathbf{\Sigma} \mathbf{V}^T$$
Trong đó:
\begin{itemize}
    \item $\mathbf{U} \in \mathbb{R}^{M \times M}$: Ma trận chứa các Left Singular Vectors (Các vector riêng của $\mathbf{A} \mathbf{A}^T$).
    \item $\mathbf{\Sigma} \in \mathbb{R}^{M \times N}$: Ma trận đường chéo chứa các giá trị đơn lẻ $\sigma_1 \ge \sigma_2 \ge \dots \ge 0$ (Đo mức độ năng lượng/thông tin).
    \item $\mathbf{V} \in \mathbb{R}^{N \times N}$: Ma trận chứa các Right Singular Vectors (Các vector riêng của $\mathbf{A}^T \mathbf{A}$).
\end{itemize}

\textbf{PCA (Principal Component Analysis):}
1. Trừ trung bình dữ liệu: $\mathbf{X}_{\text{centered}} = \mathbf{X} - \boldsymbol{\mu}$.
2. Tính ma trận hiệp phương sai: $\mathbf{S} = \frac{1}{N} \mathbf{X}_{\text{centered}}^T \mathbf{X}_{\text{centered}}$.
3. Tìm $K$ vector riêng tương ứng với $K$ giá trị riêng lớn nhất của $\mathbf{S}$.
4. Chiếu dữ liệu xuống $K$ trục thành phần chính này để nén dung lượng mà giữ được phương sai cực đại.

\newpage
% ════════════════════════════════════════════════════════════════
\section{Trụ cột 2: Giải tích Vector, Ma trận Hessian/Jacobian \& Tối ưu hóa}

\subsection{Từ điển Ký hiệu Vi tích phân \& Tối ưu}
\begin{symbolbox}[Ký hiệu Đạo hàm \& Ma trận Cấp cao]
\begin{itemize}
    \item $\nabla f(\mathbf{x}) = \left[ \frac{\partial f}{\partial x_1}, \frac{\partial f}{\partial x_2}, \dots, \frac{\partial f}{\partial x_D} \right]^T$: Vector Gradient. Chỉ hướng tăng nhanh nhất của $f$.
    \item $\mathbf{J} \in \mathbb{R}^{E \times D}$: Ma trận Jacobian chứa các đạo hàm riêng cấp 1 của hàm vector $\mathbf{f}: \mathbb{R}^D \rightarrow \mathbb{R}^E$.
    \item $\mathbf{H}_f \in \mathbb{R}^{D \times D}$: Ma trận Hessian chứa các đạo hàm riêng cấp 2 ($H_{ij} = \frac{\partial^2 f}{\partial x_i \partial x_j}$).
    \item $\mathbf{x}^* = \arg\min_{\mathbf{x}} f(\mathbf{x})$: Điểm $\mathbf{x}^*$ làm cho hàm $f(\mathbf{x})$ đạt giá trị nhỏ nhất.
    \item $\eta$ (hoặc $\gamma$): Tốc độ học (Learning Rate) trong thuật toán Gradient Descent.
\end{itemize}
\end{symbolbox}

\subsection{Ma trận Jacobian \& Ma trận Hessian}
\textbf{Ma trận Jacobian ($\mathbf{J}$):} Khi hàm có $D$ đầu vào và $E$ đầu ra:
$$\mathbf{J} = \begin{bmatrix} 
\frac{\partial f_1}{\partial x_1} & \frac{\partial f_1}{\partial x_2} & \dots & \frac{\partial f_1}{\partial x_D} \\
\frac{\partial f_2}{\partial x_1} & \frac{\partial f_2}{\partial x_2} & \dots & \frac{\partial f_2}{\partial x_D} \\
\vdots & \vdots & \ddots & \vdots \\
\frac{\partial f_E}{\partial x_1} & \frac{\partial f_E}{\partial x_2} & \dots & \frac{\partial f_E}{\partial x_D}
\end{bmatrix}$$

\textbf{Ma trận Hessian ($\mathbf{H}_f$):} Bảng đạo hàm cấp 2 đối xứng của hàm vô hướng:
$$\mathbf{H}_f = \begin{bmatrix}
\frac{\partial^2 f}{\partial x_1^2} & \frac{\partial^2 f}{\partial x_1 \partial x_2} & \dots & \frac{\partial^2 f}{\partial x_1 \partial x_D} \\
\frac{\partial^2 f}{\partial x_2 \partial x_1} & \frac{\partial^2 f}{\partial x_2^2} & \dots & \frac{\partial^2 f}{\partial x_2 \partial x_D} \\
\vdots & \vdots & \ddots & \vdots \\
\frac{\partial^2 f}{\partial x_D \partial x_1} & \frac{\partial^2 f}{\partial x_D \partial x_2} & \dots & \frac{\partial^2 f}{\partial x_D^2}
\end{bmatrix}$$

\subsection{Xác định Dấu của Ma trận \& Tìm Cực trị}
\textbf{Điều kiện Kiểm tra Dấu Ma trận Hessian tại Điểm dừng $\mathbf{x}^*$ ($\nabla f(\mathbf{x}^*) = \mathbf{0}$):}

\begin{table}[h]
\centering
\caption{Bảng Tiêu chuẩn Xác định Dấu \& Kết luận Cực trị}
\begin{tabular}{p{3.5cm}p{4cm}p{4cm}p{3.5cm}}
\toprule
\textbf{Loại Dấu Hessian} & \textbf{Điều kiện Trị riêng ($\lambda_i$)} & \textbf{Tiêu chuẩn Sylvester ($\Delta_k$)} & \textbf{Loại Cực trị} \\
\midrule
\textbf{Xác định Dương} & Tất cả $\lambda_i > 0$ & $\Delta_1 > 0, \Delta_2 > 0, \Delta_3 > 0$ & \textbf{Cực tiểu địa phương} \\
\textbf{Xác định Âm} & Tất cả $\lambda_i < 0$ & $\Delta_1 < 0, \Delta_2 > 0, \Delta_3 < 0$ (Đan dấu) & \textbf{Cực đại địa phương} \\
\textbf{Không xác định} & Có cả $\lambda_i > 0$ và $\lambda_j < 0$ & Không thỏa các tiêu chuẩn lồi/lõm & \textbf{Điểm Yên ngựa} \\
\bottomrule
\end{tabular}
\end{table}

\begin{examplebox}[Bài toán Tìm Cực trị Đầy đủ]
Tìm cực trị của hàm $f(x, y) = x^2 + y^2 - xy + x - y$.
\begin{enumerate}
    \item \textbf{Tìm Điểm dừng:} 
    $$\nabla f(x, y) = \left( \frac{\partial f}{\partial x}, \frac{\partial f}{\partial y} \right) = (2x - y + 1, 2y - x - 1)$$
    Giải hệ $\nabla f = (0, 0) \iff \begin{cases} 2x - y = -1 \\ -x + 2y = 1 \end{cases} \iff (x^*, y^*) = \left(-\frac{1}{3}, \frac{1}{3}\right)$.
    
    \item \textbf{Lập Ma trận Hessian:}
    $\frac{\partial^2 f}{\partial x^2} = 2, \quad \frac{\partial^2 f}{\partial x \partial y} = -1, \quad \frac{\partial^2 f}{\partial y^2} = 2 \implies \mathbf{H}_f = \begin{bmatrix} 2 & -1 \\ -1 & 2 \end{bmatrix}$.
    
    \item \textbf{Xét Dấu Hessian:}
    Định thức con $\Delta_1 = 2 > 0$, $\Delta_2 = (2 \times 2) - (-1)^2 = 3 > 0$.
    Hessian \textbf{Xác định dương} $\implies \left(-\frac{1}{3}, \frac{1}{3}\right)$ là điểm \textbf{Cực tiểu địa phương}.
\end{enumerate}
\end{examplebox}

\subsection{Các Thuật toán Tối ưu hóa trong AI}
\begin{itemize}
    \item \textbf{Gradient Descent (GD):} $\mathbf{w}_{k+1} = \mathbf{w}_k - \eta \nabla f(\mathbf{w}_k)$.
    \item \textbf{Stochastic Gradient Descent (SGD):} Cập nhật trọng số chỉ dựa trên 1 mẫu ngẫu nhiên $i$: $\mathbf{w}_{k+1} = \mathbf{w}_k - \eta \nabla f_i(\mathbf{w}_k)$.
    \item \textbf{Momentum:} Tích lũy vận tốc $v_{k+1} = \beta v_k + \eta \nabla f(\mathbf{w}_k)$, cập nhật $\mathbf{w}_{k+1} = \mathbf{w}_k - v_{k+1}$.
    \item \textbf{Adam Optimizer:} Thuật toán tối ưu phổ biến nhất, kết hợp cả Động lượng (Trung bình động cấp 1) và RMSProp (Trung bình động bình phương cấp 2).
\end{itemize}

\newpage
% ════════════════════════════════════════════════════════════════
\section{Trụ cột 3: Xác suất Thống kê, Phân phối \& Ước lượng Tham số}

\subsection{Từ điển Ký hiệu Thống kê}
\begin{symbolbox}[Ký hiệu Xác suất \& Phân phối]
\begin{itemize}
    \item $P(A \mid B) = \frac{P(A \cap B)}{P(B)}$: Xác suất có điều kiện của $A$ khi biết $B$.
    \item $\mathbb{E}[X]$: Kỳ vọng (Giá trị trung bình) của biến ngẫu nhiên $X$.
    \item $\text{Var}(X) = \mathbb{E}[(X - \mu)^2] = \mathbb{E}[X^2] - (\mathbb{E}[X])^2$: Phương sai (Mức độ phân tán).
    \item $\mathbf{\Sigma} = \mathbb{E}[(\mathbf{X} - \boldsymbol{\mu})(\mathbf{X} - \boldsymbol{\mu})^T]$: Ma trận Hiệp phương sai.
    \item $\mathcal{N}(\mu, \sigma^2)$: Phân phối Chuẩn (Gaussian) 1 chiều: $p(x) = \frac{1}{\sqrt{2\pi\sigma^2}} \exp\left( -\frac{(x - \mu)^2}{2\sigma^2} \right)$.
    \item $\mathcal{N}(\boldsymbol{\mu}, \mathbf{\Sigma})$: Phân phối Chuẩn nhiều chiều: $p(\mathbf{x}) = \frac{1}{(2\pi)^{D/2} |\mathbf{\Sigma}|^{1/2}} \exp\left( -\frac{1}{2} (\mathbf{x} - \boldsymbol{\mu})^T \mathbf{\Sigma}^{-1} (\mathbf{x} - \boldsymbol{\mu}) \right)$.
\end{itemize}
\end{symbolbox}

\subsection{3 Phương pháp Ước lượng Tham số Cốt lõi}

\begin{mindsetbox}[So sánh MLE, MAP và Bayesian Inference]
Khi cần tìm bộ tham số $\boldsymbol{\theta}$ cho mô hình từ tập dữ liệu $\mathcal{D} = \{x_1, x_2, \dots, x_N\}$:
\begin{enumerate}
    \item \textbf{MLE (Maximum Likelihood Estimation):} Xem $\boldsymbol{\theta}$ là số cố định chưa biết. Tìm $\boldsymbol{\theta}$ làm cho xác suất xuất hiện tập dữ liệu là lớn nhất:
    $$\hat{\boldsymbol{\theta}}_{\text{MLE}} = \arg\max_{\boldsymbol{\theta}} P(\mathcal{D} \mid \boldsymbol{\theta}) = \arg\min_{\boldsymbol{\theta}} \left[ -\sum_{i=1}^N \log P(x_i \mid \boldsymbol{\theta}) \right]$$
    
    \item \textbf{MAP (Maximum A Posteriori):} Bổ sung tri thức có sẵn ban đầu (Prior $P(\boldsymbol{\theta})$) để tránh hiện tượng quá tự tin khi ít dữ liệu:
    $$\hat{\boldsymbol{\theta}}_{\text{MAP}} = \arg\max_{\boldsymbol{\theta}} P(\boldsymbol{\theta} \mid \mathcal{D}) = \arg\max_{\boldsymbol{\theta}} \left[ \sum_{i=1}^N \log P(x_i \mid \boldsymbol{\theta}) + \log P(\boldsymbol{\theta}) \right]$$
    
    \item \textbf{Bayesian Inference:} Không chọn 1 điểm $\boldsymbol{\theta}$ duy nhất mà tính toàn bộ phân phối hậu định $P(\boldsymbol{\theta} \mid \mathcal{D})$. Dự đoán bằng cách lấy tích phân trên toàn bộ không gian tham số:
    $$P(y^* \mid x^*, \mathcal{D}) = \int P(y^* \mid x^*, \boldsymbol{\theta}) P(\boldsymbol{\theta} \mid \mathcal{D}) d\boldsymbol{\theta}$$
\end{enumerate}
\end{mindsetbox}

\begin{examplebox}[Chứng minh MLE cho Phân phối Chuẩn 1 chiều]
Cho $x_1, \dots, x_N \sim \mathcal{N}(\mu, \sigma^2)$. Tìm $\hat{\mu}_{\text{MLE}}$:
\begin{enumerate}
    \item Hàm Likelihood: $L(\mu) = \prod_{i=1}^N \frac{1}{\sqrt{2\pi\sigma^2}} \exp\left( -\frac{(x_i - \mu)^2}{2\sigma^2} \right)$.
    \item Log-Likelihood: $\ln L(\mu) = -\frac{N}{2} \ln(2\pi\sigma^2) - \frac{1}{2\sigma^2} \sum_{i=1}^N (x_i - \mu)^2$.
    \item Đạo hàm theo $\mu$ cho bằng 0: 
    $$\frac{d \ln L}{d\mu} = \frac{1}{\sigma^2} \sum_{i=1}^N (x_i - \mu) = 0 \iff \sum_{i=1}^N x_i - N\mu = 0 \iff \hat{\mu}_{\text{MLE}} = \frac{1}{N} \sum_{i=1}^N x_i$$
    \textit{Kết luận:} MLE của trung bình chính là Trung bình cộng của dữ liệu quan sát!
\end{enumerate}
\end{examplebox}

\newpage
% ════════════════════════════════════════════════════════════════
\section{Trụ cột 4: Mô hình Học máy, Hồi quy, Phân loại \& Mạng Neural}

\subsection{Hồi quy Tuyến tính (Linear Regression)}
\textbf{Mô hình:} $\hat{y} = \mathbf{w}^T \mathbf{x} + b$.

\textbf{Hàm Mất Mát (MSE Loss):} $L(\mathbf{w}) = \frac{1}{N} \|\mathbf{y} - \mathbf{X}\mathbf{w}\|_2^2$.

\textbf{Nghiệm Đóng (Normal Equation):}
$$\mathbf{w}^* = (\mathbf{X}^T \mathbf{X})^{-1} \mathbf{X}^T \mathbf{y}$$

\subsection{Hồi quy Luận lý (Logistic Regression) \& Softmax Classifier}
\textbf{Logistic Regression (Phân loại nhị phân):}
$$\hat{y} = \sigma(z) = \frac{1}{1 + e^{-z}}, \quad \text{với } z = \mathbf{w}^T \mathbf{x} + b$$
Hàm Mất mát Binary Cross-Entropy (BCE):
$$L(\mathbf{w}) = -\frac{1}{N} \sum_{i=1}^N \left[ y_i \log(\hat{y}_i) + (1 - y_i) \log(1 - \hat{y}_i) \right]$$

\textbf{Softmax Classifier (Phân loại $C > 2$ lớp):}
$$P(Y = k \mid \mathbf{x}) = \frac{e^{s_k}}{\sum_{j=1}^C e^{s_j}}, \quad \text{với } s_k = \mathbf{w}_k^T \mathbf{x} + b_k$$

\subsection{Mạng Neural \& Thuật toán Lan truyền ngược (Backpropagation)}

\begin{mindsetbox}[Quy tắc 4 Cổng Tính toán trong Pha Lùi (Backward Pass)]
Trong Đồ thị Tính toán (Computational Graph), tín hiệu Gradient truyền ngược tuân theo Chain Rule:
$$\text{Downstream Grad} = \text{Upstream Grad} \times \text{Local Grad}$$
\begin{enumerate}
    \item \textbf{Cổng Cộng (`+` Add Gate) -- Bộ Phân Phối:} Phân phối nguyên vẹn Upstream Grad cho tất cả các nhánh đầu vào.
    \item \textbf{Cổng Nhân (`*` Mul Gate) -- Bộ Tráo Đổi:} Nhân Upstream Grad với giá trị đầu vào của nhánh đối diện.
    \item \textbf{Cổng Max (`max` Max Gate) -- Bộ Định Tuyến:} Dẫn trọn vẹn Upstream Grad về nhánh có giá trị lớn hơn; nhánh nhỏ hơn nhận Grad = 0.
    \item \textbf{Cổng Phân Nhánh (Fork Gate) -- Bộ Cộng Dồn:} Cộng tổng toàn bộ Gradient gửi về từ các nhánh.
\end{enumerate}
\end{mindsetbox}

\begin{examplebox}[Tính Backpropagation từng bước qua Đồ thị Tính toán]
Xét phép tính $f(x, y) = x \cdot y + \max(x, y)$ tại $x = 2.0, y = 3.0$:
\begin{enumerate}
    \item \textbf{Forward Pass:}
    \begin{itemize}
        \item Tích $a = x \cdot y = 2.0 \times 3.0 = 6.0$.
        \item Max $b = \max(x, y) = \max(2.0, 3.0) = 3.0$.
        \item Đầu ra $f = a + b = 6.0 + 3.0 = 9.0$.
    \end{itemize}
    
    \item \textbf{Backward Pass (Khởi tạo Upstream Grad $\frac{\partial L}{\partial f} = 1.0$):}
    \begin{itemize}
        \item Cổng Cộng $f = a + b$: Phân phối nguyên vẹn $\implies \frac{\partial L}{\partial a} = 1.0$, $\frac{\partial L}{\partial b} = 1.0$.
        \item Cổng Max $b = \max(x, y)$: Vì $y = 3.0 > x = 2.0 \implies \frac{\partial b}{\partial y} = 1, \frac{\partial b}{\partial x} = 0$.
        Grad gửi về từ nhánh Max: $\frac{\partial L}{\partial x}_{\text{Max}} = 0$, $\frac{\partial L}{\partial y}_{\text{Max}} = 1.0 \times 1 = 1.0$.
        \item Cổng Nhân $a = x \cdot y$: Tráo đổi nhánh $\implies \frac{\partial a}{\partial x} = y = 3.0$, $\frac{\partial a}{\partial y} = x = 2.0$.
        Grad gửi về từ nhánh Nhân: $\frac{\partial L}{\partial x}_{\text{Mul}} = 1.0 \times 3.0 = 3.0$, $\frac{\partial L}{\partial y}_{\text{Mul}} = 1.0 \times 2.0 = 2.0$.
        \item Cổng Phân nhánh $x$ và $y$ (Cộng dồn):
        $$\frac{\partial L}{\partial x} = \frac{\partial L}{\partial x}_{\text{Mul}} + \frac{\partial L}{\partial x}_{\text{Max}} = 3.0 + 0 = \mathbf{3.0}$$
        $$\frac{\partial L}{\partial y} = \frac{\partial L}{\partial y}_{\text{Mul}} + \frac{\partial L}{\partial y}_{\text{Max}} = 2.0 + 1.0 = \mathbf{3.0}$$
    \end{itemize}
\end{enumerate}
\end{examplebox}

\newpage
% ════════════════════════════════════════════════════════════════
\section{Bảng Hệ thống Thuật toán \& Khung Truyền đạt Giảng dạy}

\subsection{Bảng Tổng hợp Thuật toán Toàn khóa}
\begin{table}[h]
\centering
\caption{Hệ thống Hóa Toàn bộ Thuật toán Môn học CS115}
\begin{tabular}{p{3.5cm}p{2.5cm}p{3.5cm}p{5.5cm}}
\toprule
\textbf{Thuật toán} & \textbf{Loại Học} & \textbf{Hàm Loss / Độ đo} & \textbf{Tư duy Cốt lõi / Ứng dụng} \\
\midrule
\textbf{Linear Regression} & Có giám sát & MSE Loss & Đường thẳng tối thiểu hóa tổng bình phương sai số. \\
\textbf{Logistic Regression} & Có giám sát & Cross-Entropy Loss & Phân loại nhị phân qua hàm Sigmoid. \\
\textbf{Softmax Classifier} & Có giám sát & Categorical Cross-Entropy & Phân loại đa lớp, tổng xác suất bằng 100\%. \\
\textbf{KNN} & Có giám sát & Khoảng cách Euclidean/$L_1$ & "Ngưu tầm ngưu": Bầu chọn theo $K$ hàng xóm gần nhất. \\
\textbf{SVM} & Có giám sát & Hinge Loss + $L_2$ Penalty & Tìm siêu phẳng tối đa hóa khoảng cách Lề (Margin). \\
\textbf{Decision Tree / Forest} & Có giám sát & Entropy / Gini Impurity & Trò chơi 20 câu hỏi Yes/No \& Bầu chọn đa số. \\
\textbf{K-Means} & Không giám sát & Inertia (Tổng khoảng cách) & Lặp 2 bước: Gán điểm vào tâm \& Tính lại vị trí tâm. \\
\textbf{PCA / SVD} & Không giám sát & Phương sai / Singular Values & Nén dữ liệu giữ lại tối đa thông tin cốt lõi. \\
\textbf{LDA} & Có giám sát & Fisher Criterion & Chiếu dữ liệu tối đa khoảng cách giữa các lớp. \\
\bottomrule
\end{tabular}
\end{table}

\subsection{Khung 3 Câu hỏi Vàng khi Giảng bài cho Người khác}
Khi truyền đạt lại bất kỳ bài toán AI nào cho bạn học, hãy luôn dẫn dắt theo đúng 3 câu hỏi chìa khóa:
\begin{enumerate}
    \item \textbf{Câu hỏi 1 (Biểu diễn):} \textit{"Bài toán này biểu diễn dữ liệu vào bằng Vector hay Ma trận nào?"}
    \item \textbf{Câu hỏi 2 (Mất mát):} \textit{"Mô hình đo độ đoán sai bằng hàm Loss nào (MSE hay Cross-Entropy)?"}
    \item \textbf{Câu hỏi 3 (Tối ưu):} \textit{"Máy tính dùng Gradient và Backpropagation ra sao để tự sửa lỗi cho đúng?"}
\end{enumerate}

\end{document}
